\uuid{siks}
\titre{Nombres premiers}
\auteur{bodin ; croizat ; sacre}
\organisation{exo7}
\licence{CC BY-NC-SA 4.0}
\datecreate{2022-03-08}
\niveau{L1}
\module{Algèbre}
\chapitre{Arithmétique dans Z}
\sousChapitre{Nombres premiers, nombres premiers entre eux}
\difficulte{3}
\theme{Nombres premiers}
\exo7qcmid{hilisit.arithmetique.039}

\contenu{
\begin{question}
\qclassification{103.04}
 Soit $p \geq 5$ un nombre premier et $N = (p+3)^2 - p^2$. Quelles sont les affirmations vraies ?
    \begin{answers}
        \bad{$2 | N$.}
        \good{$3 | N$.}
        \bad{$6 | N$.}
        \good{$p$ ne divise pas $N$.}
    \end{answers}
    \begin{explanations}
    En développant, on constate que $N = 6p+9 = 6(p+1)+3 =  3(2p+3)$. $N$ est donc un multiple de $3$, non divisible par $6$ (le reste dans la division euclidienne par $6$ est $3$). $N$ est impair (produit de deux nombres impairs) et n'est donc pas divisible par $2$.\\
    Enfin, si $p|N$, on a $p | (6p+9)$ et donc $p|9$. Puisque $p$ est premier, cela signifie que $p=3$ ce qui est impossible car $p \geq 5$.
    \end{explanations}
\end{question}
}
